% !TEX root = print.tex
% In this file you should put the actual content of the blueprint.
% It will be used both by the web and the print version.
% It should *not* include the \begin{document}
%
% If you want to split the blueprint content into several files then
% the current file can be a simple sequence of \input. Otherwise It
% can start with a \section or \chapter for instance.
\section{Root pairings}
\subsection{Definitions}
Recall that a {\it perfect pairing} is a quadruple $(R,M,N,\mathcal{L})$ where $R$ is a ring, $M$ and $N$ are $R$-modules, 
and $\mathcal{L}: M \times N \to R$ is a bilinear map such that the induced maps 
$M \to \text{Hom}_R(N,R)$ and $N \to \text{Hom}_R(M,R)$ are isomorphisms of $R$-modules. 

If $I$ is a set, denote by $\text{Perm}(I)$ the group of permutations of $I$. 

\begin{definition}\label{def:root_pairing}
    A {\it root pairing} is a tuple $(R,M,N,\mathcal{L},I,\alpha,\alpha^{\vee},s)$ where $(R,M,N,\mathcal{L})$ is a perfect pairing, 
    $I$ is a set, $\alpha: I \to M, i \mapsto \alpha_i$, $\alpha^{\vee}: I \to N, i \mapsto \alpha_i^{\vee}$, and 
    $s: I \to \text{Perm}(I), i \mapsto s_i$, with the conditions that 
\begin{enumerate}
    \item $\alpha, \alpha^{\vee}$ are injective,
    \item For all $i \in I$, $\mathcal{L}(\alpha_i, \alpha^{\vee}_i) = 2$,
    \item For all $i,j \in I$, we have 
\begin{align*}
    \alpha_{s_i(j)} &= \alpha_j - \mathcal{L}(\alpha_j, \alpha_i^{\vee})\alpha_i, \\
    \alpha_{s_i(j)}^{\vee} &= \alpha_j^{\vee} - \mathcal{L}(\alpha_i, \alpha_j^{\vee})\alpha_i^{\vee}.
\end{align*}
\end{enumerate}
\end{definition}


\subsection{Root pairing of type A}
Consider the perfect pairing $(R,M,N,\mathcal{L}) = (\mathbb{Z}, \mathbb{Z}^n, (\mathbb{Z}^n)^*, (\cdot , \cdot))$ where $(\cdot , \cdot)$ is the 
canonical pairing between $\mathbb{Z}^n$ and its algebraic dual $(\mathbb{Z}^n)^* := {\rm Hom}_{\mathbb{Z}}(\mathbb{Z}^n, \mathbb{Z})$, i.e. 
the dot product. 

We denote by $e_1, ..., e_n$ the standard basis of $\mathbb{Z}^n$ and by $e_1^*, ..., e_n^*$ the dual basis of $(\mathbb{Z}^n)^*$.

Write $[1,n] = \{1, ..., n\}$. An {\it interval} in $[1,n]$ is a subset of the form $[i,j] \stackrel{\text{denotes}}{=}\{i, i+1, ..., j\}$ 
for some $1 \leq i \leq j \leq n$. We denote by $I_n$ the set of intervals in $[1,n]$. Write $[i] := [i,i]$ for $1 \leq i \leq n$. 

A {\it signed interval} in $[1,n]$ is an element of $I_n := I_n \times \{1,-1\}. $ We denote by $I_n^+$ the set of signed intervals of the 
form $(J,1)$ and by $I_n^-$ the set of signed intervals of the form $(J,-1)$. It is clear that $I_n = I_n^+ \sqcup I_n^-$. 
Throughout the rest of this file, we write $J$ for the signed interval $(J,1)$ and $-J$ for $(J,-1)$.

Let $A = (a_{i,j})$ be the Cartan matrix of type $A_n$ defined by $a_{i,i} = 2$ for all $i \in [1,n]$, 
$a_{i,j} = -1$ if $|i-j|=1$, and $a_{i,j} = 0$ otherwise. 

Define 
\begin{align*}
    &\alpha: I_n \to \mathbb{Z}^n, & \quad \forall J \in I_n^+, \;\; \alpha_{\pm J} = \pm(\sum_{j \in J} Ae_j), \\
    &\alpha^{\vee}: I_n \to (\mathbb{Z}^n)^*, & \quad \forall J \in I_n^+, \;\; \alpha^{\vee}_{\pm J} = \pm(\sum_{j \in J} e_j^*).
\end{align*}

Finally, define $s: I_n^+ \to \text{Perm}(I_n)$ as follows. Fix $J,K \in I_n^+$ and set $L = (J \cup K) \backslash (J \cap K)$. 
It is clear that $J,K,L$ satisfy exactly one of the following four conditions.
\begin{description} 
    \item[{\bf R1.}] $L = \emptyset$.
    \item[{\bf R2.}] $L \in I_n$ and $K \subset J$, and 
    \item[{\bf R3.}] $L \in I_n$ and $K \not\subseteq J$.
    \item[{\bf R4.}] $L \not\in I_n \cup \{\emptyset\}$.
\end{description}
Then, 
\begin{equation}\label{eq:s_J_definition}
s_J(K) = \begin{cases}
    -K  & \text{in Case {\bf R1},}  \\
    -L & \text{in Case {\bf R2},}  \\
    L & \text{in Case {\bf R3}.}  \\
    K & \text{in Case {\bf R4},}  \\
\end{cases}
\end{equation}

If $J = [i,j]$ and $K = [k,l]$, then \eqref{eq:s_J_definition} can be rewritten as
\begin{equation}\label{eq:s_J_interval_definition}
            s_{[i, j]}[k, l] =
            \begin{cases}
                -[k,l]  & \text{if } (i, j) = (k, l) \\
                -[l+1,j] & \text{if } i = k, j > l \\
                [j+1, l] & \text{if } i = k, j < l \\
                -[i, k-1] & \text{if } i < k, j = l \\
                [k, i-1] & \text{if } i > k, j = l \\
                [k, j] & \text{if } i = l+1 \\
                [i, l] & \text{if } k = j+1 \\
                [k, l] & \text{otherwise}.
            \end{cases}
\end{equation}


We extend $s_J$ to $I_n$ by setting $s_{J}(-K) = -s_J(K)$ for all $K \in I_n^+$. 
We extend $s$ to $I_n$ by setting $s_{-J} = s_J$ for all $J \in I_n^+$. 

\begin{lemma}
    \label{lemma:s_J_are_permutations}
    \lean{s}
    For all $J$, we have $s_J \in {\rm Perm}(I_n)$.
\end{lemma}
\begin{proof}
    Fix $J \in I_n$. We prove that $s_J^2 = \text{id}$ by a straightforward case-by-case check using the definition of 
    $s_J$ and the four cases above. Then $s_J^{-1} = s_J$, and our claim is proved.
\end{proof}

\begin{lemma}
    \label{lemma:alpha_is_injective}
    $\alpha: I_n \to \mathbb{Z}^n$ is injective, i.e. $\alpha_{J} = \alpha_{K} \implies J = K$
\end{lemma}
\begin{proof}
    Write 
    \begin{align*}
        \alpha_J &= \varepsilon_J (\sum_{j\in +J}Ae_j), \varepsilon_J \in \{-1, 1\}\\
        \alpha_K &= \varepsilon_K (\sum_{k\in +K}Ae_k), \varepsilon_K \in \{-1, 1\}
    \end{align*}

    Then $\alpha_J = \alpha_K \implies 
    \varepsilon_J (\sum_{j\in +J}Ae_j) - \varepsilon_K (\sum_{k\in +K}Ae_k) = 0$. This reduces to
    \begin{equation*}
        \sum_{k \in +J \cup +K}c_k Ae_k = 0
    \end{equation*}
    with the following cases:
    \[ 
        c_k = 
        \begin{cases}
            \varepsilon_J & \text{if } k \in +J, k \notin +K\\
            -\varepsilon_K & \text{if } k \notin +J, k \in +K\\
            \varepsilon_J - \varepsilon_K & \text{if } k \in +J, k \in +K\\
            0 & \text{if } k \notin +J, k \notin +K
        \end{cases}
    \]

    \begin{sublemma}
        $A_n$ is invertible for all $n \in \mathbb{N}$, i.e. $\text{det}(A_n) \neq 0$
    \end{sublemma}
    \begin{proof}
        Since $A_n$ is a tridiagonal matrix with diagonal entries $= 2$ and off-diagonal entries $= -1$,
        $\text{det}(A_n) = 2\text{det}(A_{n-1})-\text{det}(A_{n-2})$. In particular, $\text{det}(A_n)
        - \text{det}(A_{n-1}) = \text{det}(A_{n-1})-\text{det}(A_{n-2})$ for all $n$. $\text{det}(A_1)=2$ 
        and $\text{det}(A_2)=3$, so $A_{n}=n+1$ for all $n$, which is clearly $\neq 0$.
    \end{proof}

    \begin{subcorollary}
        $A_n e_k$ is linearly independent $\forall k \leq n \in \mathbb{N}$. 
    \end{subcorollary}
    \begin{proof}
        It is known that a set of linearly independent vectors multiplied by an invertible matrix is still 
        linearly independent.
    \end{proof}

    By the above lemma, $Ae_k$ is linearly independent, so $\forall k\in +J \cup +K, c_k=0$. 
    Therefore, $\forall k\in +J \cup +K$, k must be in both $+J$ and $+K$. 
    In this case $\varepsilon_J - \varepsilon_K = 0$, so $\varepsilon_J = \varepsilon_K$, which means $J$ = $K$.
\end{proof}

\begin{lemma}
    \label{lemma:alpha_dual_is_injective}
    $\alpha^\vee$ is injective, i.e. $\alpha^\vee_J = \alpha^\vee_K \implies J = K$
\end{lemma}
\begin{proof}
    Clearly, the set of $e_i^*$ is linearly independent, so the same proof as Lemma \ref{lemma:alpha_is_injective} holds.
\end{proof}

\begin{lemma}
    \label{lemma:alpha_dual_alpha_is_two}
    $\forall J \in I_n, (\alpha^\vee_J, \alpha_J)=2$
\end{lemma}
\begin{proof}
    Write 
    \begin{align*}
        \alpha_J &= \varepsilon_J (\sum_{j\in +J}Ae_j), \varepsilon_J \in \{-1, 1\}\\
    \end{align*}

    Then
    \begin{align*}
        (\alpha^\vee_J, \alpha_J) &= 
        (\varepsilon_J (\sum_{j\in +J}e_j^*), \varepsilon_J (\sum_{k\in +J}Ae_j))\\
        &= \sum_{j\in +J}\sum_{k\in +J}(e_j^*, Ae_k)\\
        &= \sum_{j\in +J}\sum_{k\in +J}A_{j, k}
    \end{align*}

    For length $L= |J|$ of the interval, there are $L$ diagonal terms in the sum, 
    $2(L-1)$ pairs of diagonal-adjacent indices, each contributing $2$ and $-1$ respectively to the sum
    (by the definition of the Cartan matrix $A$ of type $A_n$).
    Therefore, $\sum_{j\in +J}\sum_{k\in +J}A_{j, k} = 2L - 2(L-1) = 2$
\end{proof}

\begin{lemma}
    \label{lemma:s_is_refl_on_alpha}
    For all $J, K \in I_n$, $\alpha_{s_J(K)}=\alpha_K - (\alpha^\vee_K, \alpha_J)\alpha_J$
\end{lemma}
\begin{proof}
    Let $J = \text{sign}(J)J$, where $J_+ = [i, j]$, $K = \text{sign}(K)K_+$, where $K_+ = [k, l]$.
    Then,
    \begin{align*}
        (\alpha^\vee_K, \alpha_J) &= \left(\sum_{u=k}^{l}e_u^*, \sum_{v=i}^{j}Ae_v\right)\\
        &= \sum_{u=k}^{l}\sum_{v=i}^{j}(e_u^*, Ae_v)\\
        &= \sum_{u=k}^{l}\sum_{v=i}^{j}A_{uv}\\
        &= 2(\text{\# of pairs }u = v) - (\text{\# of pairs }v = u+1) - (\text{\# of pairs }v = u-1)
    \end{align*}
    by the definition of $A$. Then there are three cases:
    \begin{itemize} 
        \item[{\bf C1.}] $u = v \implies v \in [i, j] \cap [k, l] \implies (\text{\# of pairs }u = v) = |[i, j] \cap [k, l]| =: L$
        \item[{\bf C2.}] $u = v+1 \implies v \in [i, j] \cap [k-1, l-1] \implies (\text{\# of pairs }u = v+1) = 
        L - \delta_{i, k} - + \delta_{i, l+1}$, and 
        \item[{\bf C3.}] $u = v-1 \implies v \in [i, j] \cap [k+1, l+1] \implies (\text{\# of pairs }u = v-1) = 
        L - \delta_{j, l} + \delta_{j, k-1}$.
    \end{itemize}
    Therefore
    \begin{align*}
        (\alpha^\vee_K, \alpha_J) &= \left(\sum_{u=k}^{l}e_u^*, \sum_{v=i}^{j}Ae_v\right)\\
        &= 2L - (L - \delta_{i, k} - \delta_{j, l} + \delta_{i, l+1}) - (L - \delta_{i, k} - \delta_{j, l} + \delta_{j, k-1})\\
        &= \delta_{i, k} + \delta_{j, l} - \delta_{i, l+1} - \delta_{j, k-1}
    \end{align*}
    In particular, the cases for $(\alpha^\vee_K, \alpha_J)$ are as follows:
    \[
    \begin{cases*}
        2 & \text{if } i = k, j = l\\
        1 & \text{if } i \neq k, j = l \lor 2 i = k, j \neq l\\
        -1 & \text{if } i = l+1 \lor k=j+1\\
        0 & \text{otherwise }
    \end{cases*}
    \]
    Applying these cases to the formula $\alpha_K - (\alpha^\vee_K, \alpha_J)\alpha_J$ (as well as the cases of $\text{sign}(J)$ and $\text{sign}(K)$), 
    it can be easily verified (case-by-case) that they are equivalent to $\alpha_{s_J(K)}$ using the definition \eqref{eq:s_J_interval_definition}.
\end{proof}

\begin{lemma}
    \label{lemma:s_is_refl_on_alpha_dual}
    For all $J, K \in I_n$, $\alpha_{s_J(K)}^\vee=\alpha^\vee_K - (\alpha^\vee_J, \alpha_K)\alpha^\vee_J$
\end{lemma}
\begin{proof}
    
\end{proof}

\begin{theorem}\label{thm:root_pairing_of_type_A}
    \uses{lemma:s_J_are_permutations,def:root_pairing}
    The 8-tuple $\left(\mathbb{Z}, \mathbb{Z}^l, \left(\mathbb{Z}^l\right)^*, (\cdot , \cdot), I_n, \alpha, \alpha^{\vee}, s\right)$ 
    is a root pairing.
\end{theorem}
\begin{proof}
\end{proof}